\documentclass[11pt]{article}

\usepackage[a4paper,margin=1in]{geometry}
\usepackage{amsmath,amssymb,amsthm}
\usepackage{mathtools}
\usepackage{hyperref}

\newtheorem{theorem}{Theorem}[section]
\newtheorem{lemma}[theorem]{Lemma}
\newtheorem{proposition}[theorem]{Proposition}
\newtheorem{corollary}[theorem]{Corollary}
\newtheorem{definition}[theorem]{Definition}
\newtheorem{remark}[theorem]{Remark}

\title{A Dyadic C2 Formal Route for Off-Critical Zeta Nonvanishing}
\author{Thiago F. M. Massensini Silva}
\date{Draft for formal review}

\begin{document}

\maketitle

\begin{abstract}
We present a dyadic C2 operator framework for studying the off-critical nonvanishing problem associated with the Riemann zeta function.  The construction separates the proof architecture into a local near-critical-zero regime, a middle-region resolvent non-cancellation regime, an edge regime, and a transfer theorem.  The formal endpoint is the official \texttt{RiemannHypothesis} proposition in mathlib.  The accompanying Lean development verifies the route from the stated C2 terminal certificate, the off-critical regional cover, and the C2/zeta identity to this endpoint.
\end{abstract}

\section{Introduction}

The purpose of this paper is to present a dyadic operator route to the Riemann zeta nonvanishing problem in the off-critical strip.

The proof architecture is organized into five layers:
\begin{enumerate}
  \item a C2 operator identity with the Riemann zeta function;
  \item a regional cover of the off-critical strip;
  \item local nonvanishing near critical-line zeros;
  \item middle-region non-cancellation via a dominance criterion;
  \item transfer to the official \texttt{RiemannHypothesis} proposition in mathlib.
\end{enumerate}

The first four layers are C2-side statements.  The final layer is a transfer
from C2-side off-critical nonvanishing to the classical zeta statement through
the formal C2/zeta identity.  The paper therefore distinguishes the analytic
certificates consumed by the C2 route from the final mathlib endpoint.

\section{Audit scope}

The review-facing Lean entry point is
\[
  \texttt{import LeanC2.PeerReview}.
\]
All paper-facing names below are exposed in the namespace
\[
  \texttt{C2.PeerReview}.
\]

The route audited here is conditional on the stated certificate packages.  The
terminal certificate consumed by the endpoint is
\[
  \texttt{RiemannHypothesisTerminalCertificate}.
\]
The bound packages, residual budgets, regional certificates, and oscillatory
witnesses exposed for audit are listed in
\[
  \texttt{docs/BOUNDS\_CERTIFICATES\_WITNESSES.md}.
\]
The theorem correspondence is maintained in
\[
  \texttt{docs/THEOREM\_MAP\_REVIEW.md}.
\]

\section{Preliminaries}

Let $s = \sigma + it \in \mathbb{C}$.  The critical line is $\sigma = 1/2$.  The off-critical strip is the part of the critical strip with $0 < \sigma < 1$ and $\sigma \ne 1/2$.

\begin{definition}[Off-critical nonvanishing]
For a complex function $F$, off-critical nonvanishing means that $F(s) \ne 0$ for every $s$ in the off-critical strip.
\end{definition}

\section{The C2 operator identity}

Let $F$ denote the review-facing C2 operator and let $c_0(s)$ denote the scalar factor appearing in the C2/zeta identity.

\begin{theorem}[Fundamental C2/zeta identity]
On the required right half-plane domain, the C2 operator satisfies
\[
  F(s) = c_0(s)\zeta(s).
\]
\end{theorem}

\begin{theorem}[Transfer theorem]
Assume that $F(s) = c_0(s)\zeta(s)$ on the off-critical strip and that $c_0(s) \ne 0$ there.  If $F$ is nonvanishing on the off-critical strip, then $\zeta$ is nonvanishing on the off-critical strip.
\end{theorem}

\section{Regional cover}

The off-critical strip is decomposed into three regimes:
\begin{enumerate}
  \item a near-critical-zero regime;
  \item a middle-region regime;
  \item an edge regime.
\end{enumerate}

\begin{theorem}[Off-critical regional cover]
If the three regional certificates cover the off-critical strip and each regional certificate proves nonvanishing of the C2 operator on its region, then the C2 operator is nonvanishing on the whole off-critical strip.
\end{theorem}

\section{Near-critical-zero regime}

The near-critical-zero regime is handled by a local punctured nonvanishing certificate.

\begin{theorem}[Near-critical-zero certificate]
Under the continuation hypotheses for the genuine C2 operator, the operator is locally nonvanishing in the punctured neighborhoods required by the near-critical-zero region.
\end{theorem}

The public Lean facade exposes this layer through:
\begin{verbatim}
GenuineNearCriticalZeroData
GenuineTaylorBoundsFamilyData
GenuineLeibnizTaylorBoundsFamilyData
\end{verbatim}

\section{Middle-region non-cancellation}

The middle-region argument is expressed as a dominance criterion.

\begin{definition}[Dominant four-level block]
The dominant four-level block is the main term extracted from the dyadic vertical decomposition.
\end{definition}

\begin{definition}[Residual budget]
The residual budget is the sum of the remaining tail, horizontal, tilt, and cutoff contributions after extraction of the main block.
\end{definition}

\begin{theorem}[Resolvent non-cancellation criterion]
If the dominant main term is strictly larger than the residual budget throughout the middle region, then the C2 operator is nonvanishing throughout the middle region.
\end{theorem}

The review-facing middle-region data include
\[
\begin{gathered}
  \texttt{ExpandedScalarMiddleRegionCertificate},\\
  \texttt{CanonicalClosedScaledResidualBudgetMiddleRegionData},\\
  \texttt{ResolventNoteMiddleRegionCertificate},\\
  \texttt{ExponentialPointwiseOscillatoryResidualUpper}.
\end{gathered}
\]
The oscillatory witness keeps the complex oscillatory main term separate from
the coarse remainder bound.

\section{Edge-region regime}

The edge-region certificate handles the remaining part of the off-critical strip not covered by the near-critical-zero and middle-region regimes.

\section{Terminal certificate}

\begin{theorem}[Terminal C2 certificate]
A terminal certificate consisting of continuation data, near-critical-zero data, middle-region data, edge-region data, and a cover proof implies off-critical nonvanishing of the C2 operator.
\end{theorem}

In the current canonical endpoint, the terminal package is represented in Lean
by
\[
  \texttt{RiemannHypothesisTerminalCertificate}.
\]
It is obtained from calibrated canonical middle-region data, with public
facade names for the surrounding inputs:
\[
\begin{gathered}
  \texttt{GenuineC2ContinuationData},\quad
  \texttt{ContinuedBulkNearCriticalZeroData},\\
  \texttt{CanonicalClosedScaledMiddleLocalData},\quad
  \texttt{CanonicalClosedScaledCoverCertificate},\\
  \texttt{OffCriticalRegionalCover}.
\end{gathered}
\]
Equivalent expanded-scalar and exact-zeta middle-region packages are exposed in
the public inventory for reviewers who want to audit the route at a more
granular level.

\begin{corollary}[Formal endpoint]
The terminal C2 certificate implies the official mathlib proposition
\[
  \texttt{RiemannHypothesis}.
\]
\end{corollary}

\section{Lean formalization}

The public Lean facade is located at
\[
  \texttt{LeanC2/PeerReview.lean}.
\]

The intended build command is
\begin{verbatim}
lake build LeanC2.Analytic.GenuineBulkConcrete LeanC2.PeerReview LeanC2
\end{verbatim}

The paper-facing theorem map is maintained in \texttt{docs/THEOREM\_MAP\_REVIEW.md}.

\subsection{Theorem correspondence}

\begin{center}
\begin{tabular}{ll}
\hline
Paper-facing statement & Lean facade declaration \\
\hline
Fundamental C2/zeta identity &
\texttt{C2FundamentalIdentityOnOffCriticalStrip} \\
Transfer to zeta nonvanishing &
\texttt{ZetaNonvanishing\_of\_C2Nonvanishing} \\
Transfer to mathlib RH &
\texttt{RiemannHypothesis\_of\_C2Nonvanishing} \\
Regional cover assembly &
\texttt{OffCriticalNonvanishing\_of\_regionalCover} \\
Middle-region criterion &
\texttt{ResolventNonCancellationCriterion\_of\_budgetBounds} \\
Terminal certificate endpoint &
\texttt{RiemannHypothesis\_of\_terminalCertificate} \\
\hline
\end{tabular}
\end{center}

\section{Reproducibility}

The reproducibility guide is maintained in \texttt{docs/REPRODUCIBILITY.md}.
The CI workflow runs the same Lean command displayed above.

\section{Discussion}

The central separation in the construction is between the C2-native nonvanishing mechanism and the final transfer to the classical zeta statement.  The formal endpoint is not a local replacement for RH, but the official mathlib \texttt{RiemannHypothesis} proposition.

\end{document}
